\subsection{Basic Definition and Examples}

\begin{definition}[Rings]
    \begin{enumerate}
        \item[(1)] A \textbf{ring} is a set together with two binary operations \( +  \) and \( \times  \) (called addition and multiplication) satisfying the following axioms: 
   \begin{enumerate}
       \item[(i)] \( (R, + )  \) is an Abelian group.
        \item[(ii)] \( \times  \) is associative: \( (a \times b) \times c  = a \times (b \times c ) \) for all \( a,b,c \in \R  \)
        \item[(iii)] The distribution law hold in \( R  \): For all \( a,b,c \in R  \)
            \begin{center}
                \( (a + b) \times  = (a \times c ) + (b \times c)  \) and \( a \times (b + c ) = (a \times b ) + (a  \times c ) \)
            \end{center}
          \item[(2)] The ring \( R  \) is \textbf{commutative} if multiplication is commutative.  
            \item[(3)] The ring \( R  \) is said to have an \textbf{identity} (or contain \( 1  \)) if there is an element \( 1 \in R  \) with 
                \[  a \times 1 = 1 \times a = a \ \ \forall a \in R.  \]
   \end{enumerate}
    \end{enumerate}
\end{definition}

\begin{definition}[Division Ring, Field]
    A ring \( R  \) with identity \( 1 \neq  0  \) is called a \textbf{division ring} if any nonzero element \( a \in R  \) has a multiplicative inverse; that is,  
    \[  \exists b \in R \ \text{such that} \ ab = ba = 1.  \]
    A commutative division ring is called a \textbf{field}.
\end{definition}

\begin{prop}
    Let \( R \) be a ring. Then 
    \begin{enumerate}
        \item[(1)] \( 0a = a 0 = 0  \) for all \( a \in R  \).
        \item[(2)] \( (-a)b = a(-b) = -(ab) \) for all \( a,b \in R  \) where \( -a  \) is the additive inverse of \( a  \)
        \item[(3)] \( (-a)(-b) = ab \) for all \( a,b \in R  \).
        \item[(4)] If \( R  \) has identity \( 1  \), then the identity is unique and \( -a = (-1) a \).
    \end{enumerate}
\end{prop}

\begin{definition}[Zero Divisors, Units]
    \begin{enumerate}
        \item[(1)] We call \( a \neq 0  \) in \( R  \) a \textbf{zero divisor} if there exists \( b \neq 0  \) in \( R  \) such that either \( ab = 0  \) or \( ba = 0  \).
        \item[(2)] Assume \( 1 \neq 0 \in R   \). We say that \( u \in R  \) is a \textbf{unit} if there exists \( v \in R  \) such that 
            \[  uv = vu = 1.  \]
        We denote the set of units of \( R  \) as \( R^{\times} \). (group of units of \( R  \))
    \end{enumerate}
\end{definition}

\begin{itemize}
    \item (2) tells us that the set of units form a group under multiplication.
    \item Zero divisors can never be units.
\end{itemize}

This result tells us that fields do not have zero divisors.

\begin{definition}[Integral Domain]
    A commutative ring with identity \( 1 \neq 0  \) is called \textbf{integral domain} if it has no zero divisors.
\end{definition}

Rings that are commutative gives rise to the cancellation property.

\begin{prop}
    Let \( R  \) be a ring and let \( a,b,c \in R  \) such that \( a  \) is not a zero divisor. Then the following are true: 
    If \( a b = ac  \), then either \( a = 0  \) or \( b = c  \). In particular, if \( a,b,c  \) are any elements in an integral domain and \( ab = ac  \), then either \( a = 0  \) or \( b = c  \).
\end{prop}

\begin{corollary}
    Any finite integral domain is a field.
\end{corollary}

Another perspective given by the cancellation law is that every nonzero element of a commutative ring that is not a zero divisor has a multiplicative inverse in a larger ring. This naturally leads us to the discussion of subrings.

\begin{definition}[Subring]
    A \textbf{subring} of the ring \( R  \) is a subgroup of \( R  \) that is closed under multiplication.
\end{definition}

To show that a subset of a ring is a subring, it is enough to show that \( S  \) is nonempty and \( S  \) is closed under multiplication and subtraction (much like the subgroup criterion in 210A).

\subsection{Examples of Rings}

\begin{definition}[Polynomials]
    \begin{enumerate}
        \item[(1)] Let \( x  \) be an indeterminate. The formal sum 
            \[  {a}_{n} x^{n} + {a}_{n-1} x^{n-1} + \cdots + {a}_{1} x + {a}_{0} \tag{\( n \geq 0  \)} \]
            with each \( {a}_{i} \in R  \) (\( 1 \leq i \leq n  \)) is called a \textbf{Polynomial} in \( x  \) with \textbf{coefficients} \( {a}_{i} \in R  \). 
        \item[(2)] If \( {a}_{n} \neq 0  \), then the polynomial is said to have \textbf{degree \( n \)}, and \( {a}_{n} x^{n} \) is called the \textbf{leading term}, and \( {a}_{n}  \) is called the \textbf{leading coefficient} (Note: the leading coefficient of the zero polynomial is taken to be zero).
        \item[(3)] We say that the polynomial is \textbf{monic} if \( {a}_{n} = 1  \).
    \end{enumerate}
\end{definition}

\begin{definition}[Polynomial Rings]
    The set of polynomials in the variable \( x  \) with coefficients in \( R  \) is called a \textbf{polynomial ring} and is denoted by \( R [x] \).
\end{definition}

Addition and multiplication defined on \( R[x] \) are the operations familiar from elementary algebra.

\begin{itemize}
    \item Addition in \( R[x] \) is defined component-wise; that is,
        \begin{align*}
        &({a}_{n}x^{n} + {a}_{n-1}x^{n-1} + \cdots + {a}_{1}x + {a}_{0}) + ({b}_{n}x^{n} + {b}_{n-1} x^{n-1} + \cdots + {b}_{1}x + {b}_{0}) \\
        &= ({a}_{n} + {b}_{n}) x^{n} + ({a}_{n-1} + {b}_{n-1}) x^{n-1} + \cdots + ({a}_{1} + {b}_{1}) x + ({a}_{0} + {b}_{0}).
        \end{align*}
    \item Multiplication is defined in the following way: 
        \begin{align*}
            &({a}_{0} + {a}_{1}x + {a}_{2} x^{2} + \cdots ) \times ({b}_{0} + {b}_{1}x + {b}_{2} x^{2} + \cdots ) \\
            &= {a}_{0} {b}_{0} + ({a}_{0} {b}_{1} + {a}_{1} {b}_{0})x + ({a}_{0} {b}_{2} + {a}_{1} {b}_{1} + {a}_{2}{b}_{0}) x^{2} + \cdots 
        \end{align*}
        In general, the coefficients of \( x^{k} \) in the product will be \( \displaystyle \sum_{ i = 0 }^{ k  } {a}_{i} {b}_{k-i}  \).
\end{itemize}

\begin{eg}[Examples of Polynomial Rings]
   \begin{itemize}
       \item \( \Z[x] \) is the polynomial ring with integer coefficients;
       \item \( \Q[x] \) is the polynomial ring with rational coefficients;
       \item \( \Z / 3\Z  [x]   \) is the polynomial ring with coefficients in \( \Z / 3 \Z  \).
   \end{itemize} 
   Here is an example computation: If \( p(x) = x^{2} + \overline{2} x + \overline{1} \) and \( q(x) = x^{3} + x + \overline{2} \), then
   \begin{align*}
       p(x) + q(x) &= x^{3} + x^{2} + (\overline{2} + \overline{1}) x + (\overline{1}+ \overline{2}) \\
                   &= x^{3} + x^{2} + \overline{3} x + \overline{3} \\
                   &= x^{3} + x^{2}.
   \end{align*}
\end{eg}

\begin{prop}[Proposition 4 Dummit and Foote]
    Let \( R  \) be an integral domain and let \( p(x) , q(x) \) be nonzero elements of \( R[x] \). Then
    \begin{enumerate}
        \item[(1)] The degree of \( p(x)q(x) =  \) the degree of \( p(x)  \) \( + \) the degree of \( q(x) \). 
        \item[(2)] The units of \( R[x] \) are just units of \( R  \).
        \item[(3)] \( R[x] \) is an integral domain.
    \end{enumerate}
\end{prop}
\begin{proof}
    \begin{enumerate}
        \item[(1)] Let \( p(x) \) and \( q(x) \) be polynomials in \( R[x] \) with leading terms \( {a}_{n}x^{n} \) and \( {b}_{m} x^{m} \), respectively. Then the leading term of \( p(x)q(x)  \) is \( {a}_{n}{b}_{n}  x^{n+m} \) with \( {a}_{n}{b}_{m} \neq 0 \). This tells us that \( \text{deg}[p(x)q(x)] = \text{deg}[p(x)] + \text{deg}[q(x)]    \).
        \item[(2)] If \( p(x) \) is a unit in \( R[x] \), then there exists a  polynomial \( q(x) \) such that \( p(x)q(x) = 1  \). This tells us that \( \text{deg} \ p(x) + \text{deg} \ q(x) = 0    \), so both \( p(x) \) and \( q(x)  \) are elements of \( R  \), and so are units in \( R  \) since their product is \( 1  \). 
        \item[(3)] If \( R  \) has no zero divisors, then neither does \( R[x] \) (because \( R \subseteq R[x] \)).
    \end{enumerate}
\end{proof}

\begin{itemize}
    \item If \( R  \) does have zero divisors, then \( R[x] \) also has them since \( R \subseteq R[x] \).
    \item If \( f(x)  \) is a zero divisor in \( R[x] \), then \(c f(x) = 0  \) for some nonzero \( c \in R  \).
    \item If \( S  \) is a subring of \( R  \), then \( S[x]  \) is a subring of \( R[x] \).
\end{itemize}

